Whenever $\left\{(S_i,\mu_i):i\in I\right\}$ is a collection of Boolean inverse monoids with invariant means and $\mathcal{U}$ is an ultrafilter of sets on $I$, the ultraproduct $\prod_\mathcal{U} S_i$ will be considered with the usual structure described in Subsection \ref{subsectionultraproductsofbooleaninversemonoids}. Moreover, if $(\G,\mu)$ is a probability measure-preserving groupoid (Definition \ref{definitionpmpgroupoid}), we will not make a distinction between $\BOR(\G,\mu)$ and $\MEAS(\G,\mu)$, unless strictly necessary.

The same way as the properties of a measure space $(X,\mu)$ are usually described in terms of its Borel sets and the measure, the dynamical properties of a probability measure-preserving groupoid $(\G,\mu)$ will be described in terms of the semigroup $\MEAS(\G,\mu)$. Since soficity is a property which describes the groupoids which can be modelled by finite groupoids, in the sense of the previous section, we start by describing the latter class.

\begin{definition}\nomenclature[Z]{$\mu_\#$}{Normalized counting measure on a finite set}\index{Normalized counting measure}\label{definitionnormalizedcountingmeasure}
If $X$ is a finite set, we denote by $\# X$ the cardinality of $X$, and we let $\mu_\#$ be the \emph{normalized counting measure} on $X$ (as a discrete Borel space): $\mu_\#(A)=\#A/\#X$ for all $A\subseteq X$.
\end{definition}

\begin{proposition}
If $\G$ is a connected finite groupoid, then the only invariant probability measure on $\G[0]$ (as a discrete Borel space) is the normalized counting measure $\mu_\#$ on $\G[0]$.
\end{proposition}
\begin{proof}
If $\mu$ is an invariant measure on $\G[0]$, then for all $x,y\in\G[0]$ we can choose $a\in\G$ with $\so(a)=x$ and $\ra(a)=y$. In particular, $A=\left\{a\right\}$ is a bisection of $\G$, so
\[\mu(\left\{x\right\})=\mu(A^*A)=\mu(AA^*)=\mu(\left\{y\right\})\]
(where $A^*=\left\{a^{-1}:a\in A\right\}$ is the inverse of $A$), and so $\mu$ is a multiple of the counting measure. If $\mu$ is a probability measure then $\mu=\mu_\#$.\qedhere
\end{proof}

In particular, if $X$ is a finite set and we consider the finest equivalence relation $X^2$ on $X$ as a (discrete) probability measure-preserving groupoid, the only invariant measure on $X$, identified as the unit space $(X^2)^{(0)}$ is the normalized counting measure $\mu_\#$ on $X$.

Whenever necessary, we will identify $\mu_\#$ with a normalized, faithful invariant mean on $\BOR(X^2)$ (as in Example \ref{exampleinvariantmeansaremeasures}) and so it induces a metric $d_\#$ on $\BOR(X^2)$, which is usually called the \emph{normalized Hamming distance}, and the associated trace $\tr_\#$ (Definition \ref{definitiontrace}).\nomenclature[Z]{$\tr_\#$}{Trace associated to the normalized counting measure}\nomenclature[Z]{$d_\#$}{Normalized Hamming distance}\index{Hammming distance}

Since $X^2$ is an equivalence relation, we identify $\BOR(X^2)=\mathcal{I}(X)$, as in Example \ref{examplebisectionsofequivalencerelationarefunctions}, in which case the trace and distance become
\[\tr_\#(f)=\mu_\#\left\{x\in\dom(f):f(x)=x\right\}\]
\[d_\#(f,g)=\mu_\#(\dom(f)\triangle\dom(g))+\mu_\#\left\{x\in\dom(f)\cap\dom(g):f(x)\neq g(x)\right\}\]

Following the notation introduced in Definition \ref{definitionmeasuredsemigroup}, we will denote this Boolean inverse semigroup endowed with this invariant mean and the associated metric and trace as $\MEAS(X^2,\mu_\#)$ or simply $\MEAS(X^2)$. We will prove that these groupoids actually serve as models for arbitrary finite probability measure-preserving groupoids.

\begin{proposition}\label{propositionfinitegroupoidsaresofic}
\begin{enumerate}
\item[(a)] If $(\G,\mu)$ is a connected finite probability measure-preserving groupoid, then there exists a finite set $X$ and an isometric semigroup embedding
\[\theta:\MEAS(\G,\mu)\to\MEAS(X^2,\mu_\#).\]
\item[(b)] If $\G$ is a finite probability measure-preserving groupoid, then there are finite sets $X_n$, $n\in\mathbb{N}$ and an isometric semigroup embedding $\Theta:\MEAS(\G)\to\prod_\mathcal{U}\MEAS(X_n^2)$, where $\mathcal{U}$ is any free ultrafilter on $\mathbb{N}$.
\end{enumerate}
\end{proposition}

Before proceeding with the proof, we introduce some notation which will be useful here as well as later in this chapter.

\begin{definition}\nomenclature[G]{$\sum_n t_n\G_n$}{Convex combination groupoid}\index{Convex combination groupoid}
If $\left\{(\G_n,\mu_n)\right\}_n$ is a sequence of probability measure-preserving groupoids and $t_n$ are non-negative numbers such that $\sum_n t_n=1$, we construct the \emph{convex combination groupoid} $\G=\sum t_n\G_n$ as follows: As groupoid, $\G$ is the coproduct $\G=\sqcup_n\G_n$, endowed with the Borel structure generated by the Borel structures of each $\G_n$, and the measure $\mu$ on $\G$ is given by $\mu(A)=\sum_n t_n\mu_n(A\cap\G_n)$.
\end{definition}

\begin{proof}[Proof of \ref{propositionfinitegroupoidsaresofic}]
\begin{enumerate}
\item[(a)] Using \ref{theoremclassificationofgroupoids}(a), suppose $\G=G\times Y^2$, where $G$ is a finite group, $Y$ is a finite set, and $Y^2$ is the coarsest equivalence relation on $Y$.

Moreover, since $\G$ is connected then the invariant probability measure on $\G[0]=Y$ is the normalized counting measure on $Y$. Let $X=G\times Y$ and set $\pi:\MEAS(G\times Y^2)\to\MEAS(X^2)$ by
\[\theta(A)=\left\{((gh,y),(h,x)):(g,(y,x))\in A,h\in G\right\}\]
which is easy to verify to be an isometric embedding.
%
\item[(b)] First suppose $\G=t\H+(1-t)\mathcal{K}$, where $\H,\mathcal{K}$ are connected finite probability measure-preserving groupoids. If $t=0$ then the natural inclusion $\BOR(\mathcal{K})\subseteq\BOR(\G)$ factors to an isometric isomorphism $\MEAS(\mathcal{K})\to\MEAS(\G)$, and similarly if $t=1$, so we can assume $t\neq 0,1$. Suppose further that $t$ is rational, say $t=p/q$ where, $p,q\in\mathbb{N}$, $0<p<q$.

By (a), take finite sets $X$ and $Y$ and isometric embeddings $\pi_{\H}:\MEAS(H)\to\MEAS(X^2)$ and $\pi_\mathcal{K}:\MEAS(\mathcal{K})\to\MEAS(Y^2)$. Let $[q]=\left\{0,1,\ldots,q-1\right\}$ be a finite set with $q$ elements, and set $Z=[q]\times X\times Y$ and $\theta:\MEAS(\G)\to\MEAS(Z^2)$ by
\begin{align*}
\theta(A)&=\left\{((j,x_2,y),(j,x_1,y)):0\leq j\leq p-1,y\in Y,(x_2,x_1)\in\pi_H(A\cap \H)\right\}\\
&\hspace{1.5cm}\cup\left\{((j,x,y_2),(j,x,y_1):p\leq j\leq q, x\in X, (y_2,y_1)\in \pi_\mathcal{K}(A\cap\mathcal{K})\right\}
\end{align*}
so that $\theta$ is an isometric embedding.

Iterating the argument above, if $\G$ is a rational convex combination of connected groupoids then $\MEAS(\G)$ embeds isometrically into some measured semigroup $\MEAS(Z^2)$ for some finite set $Z$.

Now given an arbitrary finite probability measure-preserving groupoid $(\G,\mu)$, Proposition \ref{propositioneverygroupoidisacoproductofconnected} guarantees that $\G$ is a coproduct of finite connected group\-oids, say $\G=\sqcup_{i=1}^k \G_i$. As above, we can assume $\mu(\G[0]_i)\neq 0$ for all $i$, so the map $A\mapsto \mu(A)/\mu(\G[0]_i)$ is an invariant measure on $\G_i$, thus it coincides with the normalized counting measure. This implies that $\G$ is actually the convex combination
\[\G=\sum_{i=1}^k \mu(\G[0]_i)\G_i.\]
Taking a sequence of rational tuples $(t^1_n,\ldots,t^k_n)$, $n=1,2,\ldots$ such that
\[\sum_{i=1}^k t_n^i=1\quad\text{for all }n,\quad\text{and}\quad t^i_n\overset{n\to\infty}{\longrightarrow}\mu(\G[0]_i)\quad\text{for all }i,\]
we obtain isometric semigroup morphisms $\theta_n:\MEAS(\sum_i t^i_n\G[0]_i)\to\MEAS(X_n^2)$ for certain finite sets $X_n$. Define a semigroup morphism
\[\Theta:\MEAS(\G)\to\prod_\mathcal{U}\MEAS(X_n^2),\qquad \Theta(A)=(\theta_n(A))_\mathcal{U},\]
%(where we identify $\BOR(G)=\BOR(\sum_i t^i_n\G_i)$ with $\MEAS(\sum_i t^i_n\G_i)$ for all $n$),
which is isometric with respect to $d_\mu$.\qedhere%, and so it factors to an isometric semigroup morphism of $\MEAS(\G,\mu)$, as we wanted.\qedhere
\end{enumerate}
\end{proof}